\chapter{RESULTS AND DISCUSSION}
\label{chap:results}

This chapter reports the empirical results of the data-collection pipeline,
the four core analytical modules (price tiering, LQS, BMS, sentiment), the
two evaluated predictive modules (RoBERTa sentiment classification and
Prophet price forecasting), and the rule-based alert engine. Unless otherwise
noted, descriptive statistics correspond to the most recent full daily
snapshot (\textbf{2026-04-29}), and aggregate counts span the full
collection window (\textbf{2026-04-17 to 2026-04-29, 13 consecutive days}).

\section{Descriptive Statistics}

\subsection{Data Collection Summary}

Over the 13-day collection window, the GitHub Actions pipeline successfully
collected:

\begin{itemize}
  \item \textbf{1{,}756} category ranking rows across 3 categories
  (theoretical maximum $13 \times 3 \times 50 = 1{,}950$ rows; observed
  completeness $\approx 90\%$, the gap stemming from occasional Apify actor
  rate-limiting or empty datasets).
  \item \textbf{409} product-level daily snapshots for the 30 watchlist
  ASINs ($\approx 105\%$ of the theoretical $13 \times 30 = 390$, with
  duplicates from manual reruns absorbed by the schema's day-level uniqueness
  constraint and re-resolved retries on partial failures).
  \item \textbf{3{,}986} customer reviews collected weekly across all 30
  watchlist ASINs.
  \item \textbf{1{,}060} entrant / exit events --- 528 entrants, 532 exits
  (99 of these involving Top 10 positions).
  \item \textbf{1{,}248} alerts generated by the rule engine.
  \item \textbf{0} image-change events (no listing image
  changed by more than the pHash Hamming-distance threshold of 10 bits over
  the 13-day window; the analytics path is fully wired and stored 0
  events for that period).
\end{itemize}

Table~\ref{tab:data_summary} reports the per-category descriptive snapshot
on 2026-04-29.

\begin{table}[H]
  \centering
  \begin{tabular}{lrrr}
    \toprule
    \textbf{Metric} & \textbf{Gaming KB} & \textbf{TWS} & \textbf{Portable Charger} \\
    \midrule
    Top-50 ASINs ranked       & 50      & 50      & 50      \\
    Watchlist ASINs           & 10      & 10      & 10      \\
    Price min (USD)           & 15.99   & 11.99   & 12.34   \\
    Price max (USD)           & 169.99  & 298.00  & 199.99  \\
    Price mean (USD)          & 65.31   & 86.61   & 28.85   \\
    Price median (USD)        & 59.99   & 45.98   & 25.98   \\
    Mean star rating          & 4.52    & 4.33    & 4.48    \\
    Mean reviews / product    & 4{,}578 & 21{,}921& 6{,}949 \\
    Median reviews / product  & 1{,}918 & 6{,}954 & 2{,}458 \\
    Sponsored slot share      & 4.0\%   & 4.0\%   & 4.0\%   \\
    Watchlist BSR median      & 206     & 62      & 126     \\
    \bottomrule
  \end{tabular}
  \caption{Descriptive statistics by category (snapshot of 2026-04-29)}
  \label{tab:data_summary}
\end{table}

The three categories exhibit distinctly different review-volume profiles:
True Wireless Earbuds is by far the most mature (median 6{,}954 reviews per
top-50 product, with the leader exceeding 100{,}000 reviews), while Gaming
Keyboards and Portable Chargers are an order of magnitude lower at 1{,}918
and 2{,}458 respectively. The price distributions also diverge: the TWS
category spans the widest absolute range (\$11.99--\$298.00) and is the only
category in which a meaningful share of the top 50 sit at premium price
points, whereas Portable Chargers cluster tightly around \$25.

\section{Price Tier Analysis}

K-Means clustering with $k = 3$ identified three distinct price segments in
each category. Table~\ref{tab:price_tiers} summarises the cluster
characteristics on 2026-04-29.

\begin{table}[H]
  \centering
  \begin{tabular}{llrrrr}
    \toprule
    \textbf{Category} & \textbf{Tier} & \textbf{Range (USD)} & \textbf{Mean (USD)} & \textbf{ASINs} \\
    \midrule
    Gaming KB & Entry   & 14.99--69.99   & 39.76   & 42 \\
    Gaming KB & Mid     & 54.99--154.95  & 78.43   & 27 \\
    Gaming KB & Premium & 129.99--349.99 & 201.66  & 6  \\
    \midrule
    TWS       & Entry   & 9.99--54.99    & 28.04   & 39 \\
    TWS       & Mid     & 69.99--149.00  & 101.10  & 16 \\
    TWS       & Premium & 169.95--298.00 & 230.66  & 12 \\
    \midrule
    Charger   & Entry   & 12.34--25.99   & 20.46   & 37 \\
    Charger   & Mid     & 24.98--49.99   & 31.00   & 32 \\
    Charger   & Premium & 199.99--199.99 & 199.99  & 1  \\
    \bottomrule
  \end{tabular}
  \caption{K-Means price-tier characteristics per category (2026-04-29).
  ASIN counts include sponsored slots and may exceed 50 because some ASINs
  appear at multiple ranks within a single search.}
  \label{tab:price_tiers}
\end{table}

The clusters confirm the hypothesised tier structure for two of the three
categories: Gaming Keyboards and TWS show the textbook entry / mid / premium
shape, with cluster centroids well-separated by roughly $2\times$ price
multiples. Portable Chargers degenerate into a near-binary split
(\$20 vs.\ \$31) with a single \$200 outlier; this is consistent with
the commoditised pricing structure of the category, where competition
operates on small absolute price differences rather than tier-level
positioning.

The tier boundaries are not strictly disjoint --- the upper edge of the
\textit{entry} tier overlaps with the lower edge of the \textit{mid} tier in
both Gaming Keyboards and Portable Chargers --- because the K-Means objective
minimises within-cluster variance rather than enforcing hard price gaps.
This is not a defect: it accurately reflects that two products at \$55 and
\$60 may sit in different competitive tiers depending on the rest of the
distribution.

\section{Listing Quality Score}

The Listing Quality Score (LQS) was computed for all 30 watchlist ASINs on
2026-04-29. Table~\ref{tab:lqs_summary} reports per-category aggregates and
Table~\ref{tab:lqs_extremes} the top and bottom performers.

\begin{table}[H]
  \centering
  \begin{tabular}{lrr}
    \toprule
    \textbf{Category} & \textbf{ASINs} & \textbf{Mean LQS} \\
    \midrule
    Gaming Keyboards      & 10 & 95.5 \\
    Portable Chargers     & 10 & 93.1 \\
    True Wireless Earbuds & 10 & 91.5 \\
    \midrule
    \textbf{Overall}      & 30 & \textbf{93.4} \\
    \bottomrule
  \end{tabular}
  \caption{LQS aggregates by category (2026-04-29).
  Overall median 94.2; range 79.5--96.9.}
  \label{tab:lqs_summary}
\end{table}

The compressed range (79.5--96.9 across all watchlist products) is itself a
finding: top-50 Amazon listings in mature consumer-electronics categories
have largely converged on a content-quality plateau --- complete title,
five-plus bullet points, A+ content, four-plus star average, and four-figure
review counts. The LQS therefore discriminates most strongly at the long
tail rather than the head of the distribution.

\begin{table}[H]
  \centering
  \begin{tabular}{lllr}
    \toprule
    \textbf{Rank} & \textbf{ASIN} & \textbf{Category} & \textbf{LQS} \\
    \midrule
    \multicolumn{4}{l}{\textit{Top 5 watchlist ASINs by LQS}} \\
    1 & B0C9ZJHQHM & Gaming Keyboards      & 96.9 \\
    2 & B0CRTR3PMF & True Wireless Earbuds & 96.8 \\
    3 & B0CDX5XGLK & Gaming Keyboards      & 96.7 \\
    4 & B0CLLHSWRL & Gaming Keyboards      & 96.7 \\
    5 & B0CB1FW5FC & Portable Chargers     & 96.6 \\
    \midrule
    \multicolumn{4}{l}{\textit{Bottom 5 watchlist ASINs by LQS}} \\
    1 & B0FQFB8FMG & True Wireless Earbuds & 79.5 \\
    2 & B0DGHMNQ5Z & True Wireless Earbuds & 81.7 \\
    3 & B0G1PJLWLZ & True Wireless Earbuds & 88.3 \\
    4 & B0DCZ56QNL & Portable Chargers     & 90.2 \\
    5 & B0CFHHCMNS & Portable Chargers     & 91.4 \\
    \bottomrule
  \end{tabular}
  \caption{LQS top and bottom performers (2026-04-29)}
  \label{tab:lqs_extremes}
\end{table}

The bottom-performing ASINs in the TWS category (Apple AirPods Pro 3,
Apple AirPods 4, Samsung Galaxy Buds 4 Pro) are all very recently launched
flagship products. Their relatively low LQS reflects modest cumulative
review counts, which the LQS weights heavily, rather than weak listing
content per se. This is a known characteristic of the metric and does not
imply that those listings are objectively poor.

\section{Brand Momentum}

Table~\ref{tab:bms} presents the top 10 watchlist ASINs by Brand Momentum
Score on 2026-04-29.

\begin{table}[H]
  \centering
  \begin{tabular}{rllrrrr}
    \toprule
    \textbf{Rk} & \textbf{ASIN} & \textbf{Category} & \textbf{BSR vel.}
      & \textbf{Rev. vel.} & \textbf{Sent.} & \textbf{BMS} \\
    \midrule
    1  & B0CB1FW5FC & Charger & $+$0.125  & 116.0 & 0.710 & 0.752 \\
    2  & B0CRTR3PMF & TWS     & $+$0.015  & 142.0 & 0.849 & 0.739 \\
    3  & B09DT48V16 & TWS     & $+$0.087  & 59.0  & 0.629 & 0.735 \\
    4  & B0DGHMNQ5Z & TWS     & ---       & 73.0  & 0.654 & 0.715 \\
    5  & B0DD41G2NZ & TWS     & $+$0.186  & 42.0  & 0.613 & 0.710 \\
    6  & B0BQPNMXQV & TWS     & 0.000     & 64.0  & 0.588 & 0.709 \\
    7  & B07XVCP7F5 & Gaming  & $-$0.054  & 274.0 & 0.605 & 0.697 \\
    8  & B0FQFB8FMG & TWS     & ---       & 78.0  & 0.301 & 0.680 \\
    9  & B0BTYCRJSS & TWS     & $-$0.179  & 183.0 & 0.704 & 0.676 \\
    10 & B0CY2JJ4WS & Charger & 0.000     & 40.0  & 0.145 & 0.605 \\
    \bottomrule
  \end{tabular}
  \caption{Top 10 watchlist ASINs by Brand Momentum Score (2026-04-29).
  ``BSR vel.'' is daily fractional BSR improvement (positive = rank
  improving), ``Rev.\ vel.'' is the daily new review count, ``Sent.''\ is
  the latest avg sentiment score in $[-1, 1]$. Dashes indicate
  null BSR (newly launched ASINs without an established BSR window).}
  \label{tab:bms}
\end{table}

True Wireless Earbuds dominate the leaderboard (7 of the top 10 entries),
which is consistent with the category's substantially higher absolute review
volume --- the BMS review-velocity term is unbounded and therefore rewards
high-traffic categories. Two interpretive notes follow:

\begin{itemize}
  \item B0CRTR3PMF (Soundcore P30i) achieves the second-highest BMS while
  combining all three signals in a balanced way: positive BSR velocity, high
  daily review volume, and the highest sentiment score on the leaderboard
  (0.849). This is the cleanest example of organic competitive momentum in
  the watchlist.
  \item B0BTYCRJSS (Soundcore P20i, the category \#1) shows
  \emph{negative} BSR velocity ($-0.179$) but still ranks 9th on BMS due to
  high review velocity (183 / day) and strong sentiment (0.704); the
  combination is consistent with a defending market leader experiencing
  routine intra-week rank churn rather than a structural decline.
\end{itemize}

\section{Sentiment Analysis}

\subsection{Distribution by Category}

Table~\ref{tab:sentiment_cat} aggregates per-ASIN sentiment scores from
\texttt{review\_sentiment\_daily} across the 30 watchlist ASINs on
2026-04-29. The score $\bar{s}_{\text{ASIN}} \in [-1, 1]$ is the
confidence-weighted polarity defined in the methodology chapter; a value of
0.5 corresponds roughly to ``75\% positive reviews / 25\% negative''.

\begin{table}[H]
  \centering
  \begin{tabular}{lrrrr}
    \toprule
    \textbf{Category} & \textbf{ASINs} & \textbf{Avg score} & \textbf{Pos.\ ratio} & \textbf{Neg.\ ratio} \\
    \midrule
    True Wireless Earbuds & 10 & 0.564 & 77.3\% & 17.5\% \\
    Gaming Keyboards      & 10 & 0.490 & 71.3\% & 21.1\% \\
    Portable Chargers     &  5 & 0.294 & 47.0\% & 16.5\% \\
    \bottomrule
  \end{tabular}
  \caption{Per-category review sentiment (2026-04-29). Portable Chargers
  shows fewer ASINs because not every watchlist product had reviews
  re-aggregated on that exact day; the missing rows are populated on the
  next weekly review-ingest cycle.}
  \label{tab:sentiment_cat}
\end{table}

The clear ordering --- TWS (0.564) > Gaming KB (0.490) > Charger (0.294) ---
mirrors the star-rating mean from Table~\ref{tab:data_summary} only at the
extremes (Charger has the lowest sentiment but the second-highest stars at
4.48), which is itself an interesting finding: a 4+ star average can mask
substantial customer concern in the free-text content, since 3-star
``it's fine, but\ldots'' reviews and qualified 4-star reviews are
sentiment-neutral or negative even when the numeric rating is high. This
is precisely the gap that justifies adding NLP-based sentiment analysis
on top of star-rating tracking.

\subsection{Validation: Distant-Supervision Evaluation}

Applying the distant-supervision protocol introduced in the methodology
chapter to all \textbf{3{,}851} reviews collected during the window for
which both a predicted label and a star rating were available, the RoBERTa
classifier achieved:

\begin{itemize}
  \item Accuracy: \textbf{81.2\%}
  \item Macro-averaged F\textsubscript{1}: \textbf{0.594}
  \item Per-class F\textsubscript{1}: positive \textbf{0.908}, negative
  \textbf{0.715}, neutral \textbf{0.158}
\end{itemize}

The detailed per-class breakdown is reported in
Table~\ref{tab:sentiment_eval} and the confusion matrix in
Table~\ref{tab:sentiment_confusion}.

\begin{table}[H]
  \centering
  \begin{tabular}{lrrr}
    \toprule
    \textbf{Class} & \textbf{Precision} & \textbf{Recall} & \textbf{F\textsubscript{1}} \\
    \midrule
    Positive & 0.929 & 0.888 & 0.908 \\
    Negative & 0.658 & 0.782 & 0.715 \\
    Neutral  & 0.159 & 0.157 & 0.158 \\
    \bottomrule
  \end{tabular}
  \caption{RoBERTa per-class precision, recall, F\textsubscript{1} (n = 3{,}851)}
  \label{tab:sentiment_eval}
\end{table}

\begin{table}[H]
  \centering
  \begin{tabular}{lrrr}
    \toprule
    \textbf{Truth $\downarrow$ / Pred $\rightarrow$} & \textbf{Negative} & \textbf{Neutral} & \textbf{Positive} \\
    \midrule
    Negative ($n=698$)  & 546 & 76  & 76   \\
    Neutral  ($n=299$)  & 136 & 47  & 116  \\
    Positive ($n=2854$) & 148 & 173 & 2{,}533 \\
    \bottomrule
  \end{tabular}
  \caption{Confusion matrix --- RoBERTa predictions vs.\ star-rating labels}
  \label{tab:sentiment_confusion}
\end{table}

The reported metrics support a nuanced reading. Performance on the two
\emph{informative} extremes (positive F\textsubscript{1} = 0.908, negative
F\textsubscript{1} = 0.715) is solid for an off-the-shelf model applied
zero-shot to a new domain. The collapse on the neutral class
(F\textsubscript{1} = 0.158) is consistent with the literature: 3-star
reviews are an inherently ambiguous category when used as a sentiment
ground truth, since reviewers regularly assign 3 stars to text that the
model --- and most human readers --- would classify as either mildly
positive or mildly negative. Of the 299 star-3 reviews, only 47 are
predicted neutral, while 136 are predicted negative and 116 positive,
suggesting that the truth labels themselves split that way and that the
model is not systematically biased.

For the downstream BMS and LQS metrics, what matters is the polarity signal
$\bar{s}_{\text{ASIN}}$ rather than per-review classification, and the
positive/negative discrimination quality reported above is sufficient for
that purpose.

\section{Price Forecasting Backtest}

Applying the walk-forward backtest defined in the methodology chapter to
the 17 watchlist ASINs with at least 7 days of price history yielded the
pooled metrics reported in Table~\ref{tab:forecast_eval}.

\begin{table}[H]
  \centering
  \begin{tabular}{lr}
    \toprule
    \textbf{Metric} & \textbf{Value} \\
    \midrule
    ASINs evaluated                       & 17       \\
    Forecast horizon (days)               & 3        \\
    Minimum training history (days)       & 7        \\
    Pooled MAE (USD)                      & 0.601    \\
    Pooled MAPE                           & 1.16\%   \\
    Naive last-value baseline MAE (USD)   & 0.601    \\
    80\% prediction interval coverage     & 90.2\%   \\
    \bottomrule
  \end{tabular}
  \caption{Walk-forward Prophet backtest (3-day horizon, 17 watchlist ASINs)}
  \label{tab:forecast_eval}
\end{table}

Three observations are worth highlighting:

\begin{enumerate}
  \item The pooled MAE of \$0.60 corresponds to a sub-2\% MAPE, which is
  small in absolute terms but \textbf{matched exactly by the naive last-value
  baseline}. Of the 17 evaluated ASINs, 14 had completely flat
  prices in their training window (MAE = MAPE = 0 for both Prophet and
  naive); the pooled error is therefore driven by 3 ASINs with mild price
  variation (B07QGHK6Q8 MAPE 8.19\%, B07XVCP7F5 MAPE 10.81\%, B0DGHMNQ5Z
  MAPE 0.66\%). Prophet does not beat the naive baseline on a 3-day
  horizon \emph{in this dataset}, which is honest --- and expected --- given
  that two-week-old price series in stable consumer-electronics categories
  contain very little exploitable trend or seasonality.
  \item The 80\% prediction interval covers 90.2\% of hold-out actuals.
  Prophet's PIs are therefore conservatively wide, which is appropriate
  for an operational forecast that downstream skills should consume with
  uncertainty.
  \item Despite the MAE parity with the naive baseline, the Prophet pipeline
  delivers two non-trivial benefits over a hand-coded last-value rule:
  (a) it provides calibrated PIs that the alert engine can use to flag
  ``price likely to drop below \$X with $\geq$\,80\% confidence'', and (b)
  the same fitted model produces a 7-day --- not just 3-day --- forecast
  that the production \texttt{query\_price\_forecast} skill exposes to
  the Strategist agent.
\end{enumerate}

The forecaster therefore contributes \textbf{methodological scaffolding}
rather than measurable point-forecast accuracy gains under the present 13-day
data regime. Section~\ref{sec:limitations} returns to this point.

\section{Entrant and Exit Patterns}

Across the 13-day window, \textbf{1{,}060} entrant or exit events were
detected: 528 entrants and 532 exits, of which 99 (9.3\%) involved a Top 10
position. Table~\ref{tab:entrants_summary} summarises the events by
category.

\begin{table}[H]
  \centering
  \begin{tabular}{lrrr}
    \toprule
    \textbf{Category} & \textbf{Entrants} & \textbf{Exits} & \textbf{Total} \\
    \midrule
    Portable Chargers     & 204 & 201 & 405 \\
    Gaming Keyboards      & 171 & 183 & 354 \\
    True Wireless Earbuds & 153 & 148 & 301 \\
    \midrule
    \textbf{All}          & \textbf{528} & \textbf{532} & \textbf{1{,}060} \\
    \bottomrule
  \end{tabular}
  \caption{Entrant / exit events by category over 13 days}
  \label{tab:entrants_summary}
\end{table}

Portable Chargers exhibits the highest churn (405 events / $50 \times 13 =
650$ ASIN-days $\approx 62\%$ daily turnover at the rank level), consistent
with the commoditised nature of the category. True Wireless Earbuds is the
most stable, with the entrants and exits roughly balanced --- the top of
this category is dominated by entrenched brands (Apple, Soundcore, JBL)
whose ranks vary inside the Top 50 but rarely drop out of it.

\section{Alert Summary}

The rule-based alert engine produced \textbf{1{,}248} alerts across the
13-day collection window. Table~\ref{tab:alert_summary} reports the breakdown
by type and severity.

\begin{table}[H]
  \centering
  \begin{tabular}{lrrrr}
    \toprule
    \textbf{Alert Type} & \textbf{Count} & \textbf{High} & \textbf{Medium} & \textbf{Low} \\
    \midrule
    \texttt{exit}        & 648    & 0   & 648 & 0   \\
    \texttt{new\_entrant} & 593   & 51  & 0   & 542 \\
    \texttt{price\_drop} & 6      & 3   & 3   & 0   \\
    \texttt{stockout}    & 1      & 1   & 0   & 0   \\
    \texttt{image\_change} & 0    & 0   & 0   & 0   \\
    \texttt{bsr\_drop}   & 0      & 0   & 0   & 0   \\
    \midrule
    \textbf{Total}       & \textbf{1{,}248} & \textbf{55} & \textbf{651} & \textbf{542} \\
    \bottomrule
  \end{tabular}
  \caption{Alert summary by type and severity (full window)}
  \label{tab:alert_summary}
\end{table}

The dominant alert types are entrant and exit events ($\approx$ 99\% of all
alerts), which reflects the high background turnover in the category top-50
ranks. The 55 high-severity alerts (51 Top-10 entrants, 3 large price drops,
1 stockout) are the operationally meaningful subset and form the
``Critical'' section of the Strategist's daily Telegram brief.

\texttt{image\_change} alerts and \texttt{bsr\_drop} alerts both fired zero
times during the window. The image-change null result is consistent with
the descriptive finding above: top-50 sellers in mature categories do not
swap their hero image during a two-week window. The \texttt{bsr\_drop} null
result reflects the threshold (BSR worsening by $\geq 500$ positions) being
calibrated for clearly anomalous moves rather than routine intraday churn;
in production with months of data, this threshold should be re-tuned per
category.

\section{Discussion}

\subsection{Competitive Dynamics Across Categories}

The empirical results reveal three distinct competitive regimes:

\begin{enumerate}
  \item \textbf{Mature, rank-stable, high-trust} (TWS): few new entrants in
  the top tier, high review volume, high LQS dispersion driven by recent
  flagship launches that have not yet accumulated reviews.
  \item \textbf{Mid-mature, brand-differentiated} (Gaming Keyboards):
  meaningful price-tier separation, mid review volumes, the lowest sponsored
  share, and the cleanest BMS interpretation.
  \item \textbf{Commoditised, high-churn} (Portable Chargers): degenerate
  price tiers, the highest entrant/exit churn, and noticeably lower review
  sentiment despite competitive star ratings --- exactly the regime where
  free-text sentiment analysis adds the most diagnostic value over star
  averages.
\end{enumerate}

\subsection{Methodology Validation}

The pipeline achieved $\approx$ 90\% data-collection completeness on
category rankings, $\sim$ 100\% on watchlist daily snapshots, and 100\% on
weekly reviews. The two evaluated predictive components (RoBERTa sentiment,
Prophet forecast) were benchmarked under defensible protocols (distant
supervision against stars; walk-forward backtest with naive baseline), and
both were found fit for their intended downstream role: sentiment supplies
a directional polarity signal to the BMS and LQS scores rather than a
calibrated probability, and the forecast supplies a probabilistic envelope
for trend communication rather than a precise point prediction.

The architectural validation came from the multi-agent layer. All
13 skills were registered with the OpenClaw Gateway, all three Telegram
bots were paired and bound 1-to-1 to their respective agents, and the
\texttt{daily-brief-strategist} cron job was verified end-to-end --- a
fully autonomous brief was delivered to the Strategist's Telegram channel
based on the previous 24 hours of Supabase data, with no manual
intervention.

\subsection{Limitations}
\label{sec:limitations}

The most significant limitations of these results, in addition to those
already discussed at the chapter level:

\begin{itemize}
  \item \textbf{Observation window.} Thirteen days is too short for
  meaningful seasonality, holiday-effect, or longitudinal-trend analysis.
  The Prophet forecast in particular cannot exploit weekly or monthly
  patterns and falls back to a linear trend with very wide PIs.
  \item \textbf{Category coverage.} Three consumer-electronics
  sub-categories may not generalise to broader marketplace dynamics. The
  commoditisation finding for Portable Chargers, in particular, is
  category-specific.
  \item \textbf{Distant supervision.} The 81.2\% accuracy figure for the
  RoBERTa pipeline is bounded above by the noisiness of star ratings as a
  sentiment proxy; a small gold-labelled review subset would tighten the
  estimate.
  \item \textbf{Telegram polling stability.} The OpenClaw Gateway runs in a
  VMware Ubuntu VM whose NAT layer occasionally drops outbound long-poll
  requests to Telegram. Recovery is automatic but can introduce 30--60s
  message lags. A managed cloud-VM deployment would remove this artefact.
\end{itemize}
